\chapter{Calculus II --- Many Variables}
\label{ch:calculus2}

\begin{tcolorbox}[feynbox={Sky}{BBg},title={Before and After}]
\textbf{Before this chapter you need:} derivatives and the rules for computing
them (\S\ref{sec:derivative}--\S\ref{sec:rules}), second derivatives and
concavity (\S\ref{sec:secondderiv}), composition of functions
(\S\ref{sec:composition}).\\
\textbf{After it you will understand:} what a partial derivative is, why the
gradient points uphill, how the chain rule propagates through a deep composition
--- which is backpropagation --- and what $\norm{\nabla D}_2 = 1$ is asking for
in the WGAN gradient penalty.
\end{tcolorbox}

This is the load-bearing chapter of the book. Gradient descent, backpropagation
and the gradient penalty are all direct applications of what follows. Everything
here is the one-variable calculus of Chapter~\ref{ch:calculus1} applied one
variable at a time, and then assembled.

A note on ordering: this chapter needs to collect several numbers into a single
object and occasionally measure its length. It does so using only school
geometry --- Pythagoras --- and defers the general theory of such objects to
Chapter~\ref{ch:linalg}, which gives them their proper treatment.

%% ─────────────────────────────────────────────────────────────────────────────
\section{Functions of Several Variables}
\label{sec:multivar}

\situation{
The temperature in a room depends on where you stand. Move toward the radiator
and it rises; move toward the window and it falls. To state the temperature you
must give two numbers --- how far along, how far across --- and you get back one
number.
}

\problem{
A curve drawn on a page cannot represent this. Two inputs and one output need
three dimensions, and the objects we actually optimise in this thesis have
hundreds of thousands of inputs. Any intuition tied to pictures will fail; what
survives is the arithmetic.
}

\workedarith{
Take $f(x,y) = x^2 + y^2$ --- squared distance from the origin. Evaluate on a grid:

\medskip
\begin{tabular}{@{}lllll@{}}
\toprule
& $y=0$ & $y=1$ & $y=2$ & $y=3$ \\
\midrule
$x=0$ & $0$ & $1$ & $4$  & $9$  \\
$x=1$ & $1$ & $2$ & $5$  & $10$ \\
$x=2$ & $4$ & $5$ & $8$  & $13$ \\
$x=3$ & $9$ & $10$& $13$ & $18$ \\
\bottomrule
\end{tabular}

\medskip
Read across row $x=2$: values $4, 5, 8, 13$, rising as $y$ grows. Read down
column $y=1$: values $1, 2, 5, 10$, rising as $x$ grows. The single lowest entry
is $0$ at the corner $(0,0)$ --- the bottom of a bowl, now in two input
directions rather than one.

\medskip
Notice the entries equal to $5$: at $(1,2)$ and at $(2,1)$. Different points, same
height. Joining all points of equal height gives a \textbf{contour}, exactly as on
a map. Here the contours are circles about the origin.
}

\formaldef{
A function of several variables assigns one output number to each combination of
inputs, written $f : \R^n \to \R$. The notation $\R^n$ means ordered lists of $n$
real numbers; $\R^2$ is pairs $(x,y)$, $\R^3$ is triples, and so on.

A \textbf{level set} or \textbf{contour} is the collection of inputs sharing one
output value: $\{(x,y) : f(x,y) = c\}$.
}

\analogybreak{
Two inputs can be pictured as a landscape; three already cannot, and the loss
function of the TimeGAN generator in this project has $11{,}400$ inputs --- one per
weight. Nothing visual survives. What does survive is that each input can be
varied on its own while the rest are held fixed, and that is the entire content of
the next section.
}

\whyhere{
A neural network's loss is exactly such a function: inputs are the weights,
output is a single number measuring how badly the network is doing. Training
searches that space for a low point.
}

%% ─────────────────────────────────────────────────────────────────────────────
\section{Partial Derivatives}
\label{sec:partials}

\situation{
You are on a hillside at a crossroads. One road runs due east, another due north.
Standing still, you can ask two separate questions: if I step east, do I go up or
down, and how steeply? And the same for north. Two roads, two answers, one place.
}

\problem{
``The slope'' of a surface is not a single number, because the answer depends on
which way you face. But if we agree to ask only about the coordinate directions
--- one at a time, everything else frozen --- each question reduces to the
one-variable derivative already available.
}

\workedarith{
Take $f(x,y) = x^2 + y^2$ at the point $(3, 4)$, where $f = 9 + 16 = 25$.

\medskip
\textbf{Step east ($x$ changes, $y$ frozen at 4).} The function becomes
$g(x) = x^2 + 16$, a one-variable function. Its derivative is $2x$, which at
$x=3$ is $6$.

Check numerically: $f(3.001, 4) = 9.006001 + 16 = 25.006001$, so
$(25.006001 - 25)/0.001 = 6.001$, tending to $6$. \checkmark

\medskip
\textbf{Step north ($y$ changes, $x$ frozen at 3).} Now
$h(y) = 9 + y^2$, derivative $2y$, which at $y=4$ is $8$.

Check: $f(3, 4.001) = 9 + 16.008001 = 25.008001$, giving
$(25.008001-25)/0.001 = 8.001 \to 8$. \checkmark

\medskip
\textbf{The recipe.} To differentiate with respect to one variable, treat every
other variable as a constant. For $f(x,y) = x^2 y^3$:
\[
  \frac{\partial f}{\partial x} = 2x y^3 \quad(\text{$y^3$ is just a number}),
  \qquad
  \frac{\partial f}{\partial y} = 3 x^2 y^2 \quad(\text{$x^2$ is just a number}).
\]
At $(2,1)$: $\partial f/\partial x = 2 \cdot 2 \cdot 1 = 4$ and
$\partial f/\partial y = 3 \cdot 4 \cdot 1 = 12$.
}

\formaldef{
The \textbf{partial derivative} of $f$ with respect to $x_i$ is
\[
  \frac{\partial f}{\partial x_i}
  = \lim_{h \to 0} \frac{f(x_1,\ldots,x_i + h,\ldots,x_n) - f(x_1,\ldots,x_n)}{h} ,
\]
the ordinary derivative in the $i$-th coordinate with all others held fixed. The
curved $\partial$ distinguishes it from the $d$ of one-variable calculus and is a
reminder that other variables exist and are being held still.
}

\analogybreak{
Partial derivatives only ever look along the coordinate axes. A ridge running
diagonally can have zero slope due east and zero due north while dropping steeply
to the southwest --- both partials vanish and the point is still not a minimum.
The coordinate directions are a convenient basis, not a complete description, and
\S\ref{sec:gradient} is what repairs the gap.
}

\whyhere{
$\partial(\text{loss})/\partial w$ for a single weight $w$ is what an optimiser
consumes. PyTorch's \texttt{backward()} computes exactly this for every weight in
the network in one sweep.
}

%% ─────────────────────────────────────────────────────────────────────────────
\section{The Gradient}
\label{sec:gradient}

\situation{
Back at the crossroads. You know the eastward slope is $6$ and the northward
slope is $8$. You are not obliged to walk along a road --- you may set off in any
direction. Which one climbs fastest, and how steep is it?
}

\problem{
The partials give two numbers; the question asks for a direction and a
steepness. Something must convert the pair into that, and the conversion had
better be the one that actually identifies the steepest route, since gradient
descent depends on going exactly opposite to it.
}

\workedarith{
Collect the partials into an ordered pair, written in brackets:
\[
  \nabla f(3,4) = \bigl[\,6,\; 8\,\bigr] .
\]

\textbf{Claim: this pair points in the steepest uphill direction, and its length
is the steepest slope.}

\medskip
\textbf{Its length,} by Pythagoras:
\[
  \sqrt{6^2 + 8^2} = \sqrt{36 + 64} = \sqrt{100} = 10 .
\]

\medskip
\textbf{Check that no direction beats 10.} Moving a small distance $h$ in a
direction that is a fraction $a$ east and $b$ north (with $a^2 + b^2 = 1$)
changes $f$ by approximately $h(6a + 8b)$. Try several:

\medskip
\begin{tabular}{@{}llll@{}}
\toprule
Direction $(a,b)$ & $6a + 8b$ \\
\midrule
Due east $(1,0)$              & $6$ \\
Due north $(0,1)$             & $8$ \\
Diagonal $(0.707, 0.707)$     & $4.24 + 5.66 = 9.90$ \\
$(0.6, 0.8)$                  & $3.6 + 6.4 = 10.0$ \\
$(0.5, 0.866)$                & $3.0 + 6.93 = 9.93$ \\
\bottomrule
\end{tabular}

\medskip
The winner is $(0.6, 0.8)$, giving exactly $10$ --- and $(0.6, 0.8)$ is
$[6,8]$ divided by its length $10$. The gradient direction wins, and the best
achievable rate equals the gradient's length. \checkmark

\medskip
\textbf{Downhill} is the exact reverse, $[-6,-8]$, dropping at rate $10$. That is
the direction gradient descent steps in.
}

\formaldef{
The \textbf{gradient} of $f : \R^n \to \R$ is the ordered list of its partial
derivatives:
\[
  \nabla f = \left[\frac{\partial f}{\partial x_1},\;
                   \frac{\partial f}{\partial x_2},\; \ldots,\;
                   \frac{\partial f}{\partial x_n}\right] .
\]
Its \textbf{length} (Chapter~\ref{ch:linalg} calls this the Euclidean or
$2$-norm) is
\[
  \norm{\nabla f}_2 = \sqrt{\left(\frac{\partial f}{\partial x_1}\right)^2 + \cdots
                            + \left(\frac{\partial f}{\partial x_n}\right)^2} .
\]
Two facts, both verified numerically above: $\nabla f$ points in the direction of
steepest increase, and $\norm{\nabla f}_2$ is the rate of increase in that
direction. Where $f$ has a local maximum or minimum in the interior of its domain,
$\nabla f = \mathbf{0}$ --- every partial vanishes at once.
}

\analogybreak{
Steepest \emph{here} is not fastest \emph{overall}. The gradient is purely local:
it describes an infinitesimal neighbourhood and knows nothing of the terrain
beyond. Repeatedly following it can zig-zag down a narrow valley, taking many
small steps where one well-chosen direction would have done. It is greedy, and
greedy is not optimal --- which is why Chapter~10 (Optimisation)'s momentum
and adaptive methods exist at all.
}

\whyhere{
Two direct uses. First, gradient descent: every model here is trained by
repeatedly stepping along $-\nabla(\text{loss})$. Second, and more specifically,
the WGAN-GP gradient penalty is a statement about $\norm{\nabla D}_2$ --- it adds a
cost whenever the length of the critic's gradient departs from $1$. With that
definition in hand, $\norm{\nabla D}_2 = 1$ now reads concretely: whichever
direction you perturb the input, the critic's output should change at a rate of
exactly one unit per unit of perturbation. Chapter~\ref{ch:gans} explains why
that is the right thing to ask for.
}

%% ─────────────────────────────────────────────────────────────────────────────
\section{The Chain Rule and Backpropagation}
\label{sec:chainmulti}

\situation{
The factory line of \S\ref{sec:composition}, now with a quality inspector at the
end who scores each finished item. A dial on station one is nudged. The painted
part changes a little; the boxed part changes as a result; the score changes as a
result of that. How much did the score move per unit of dial?
}

\problem{
The dial and the score are separated by every station in between. Measuring the
effect by physically turning the dial and re-running the line is possible for one
dial, but a network has hundreds of thousands, and re-running once per dial is
hopelessly expensive. We need the whole answer from one pass.
}

\workedarith{
\textbf{A chain of three.} Let
\[
  u = g(x) = x + 3, \qquad v = f(u) = u^2, \qquad L = k(v) = 2v .
\]
At $x = 1$: $u = 4$, $v = 16$, $L = 32$.

Local rates: $du/dx = 1$, $dv/du = 2u = 8$, $dL/dv = 2$.

Chain them by multiplying:
\[
  \frac{dL}{dx} = \frac{dL}{dv}\cdot\frac{dv}{du}\cdot\frac{du}{dx}
  = 2 \times 8 \times 1 = 16 .
\]

\textbf{Verify by substitution.} $L = 2(x+3)^2$, so $dL/dx = 4(x+3)$, which at
$x=1$ is $16$. \checkmark

\textbf{Verify numerically.} At $x = 1.001$: $u = 4.001$, $v = 16.008001$,
$L = 32.016002$. Then $(32.016002 - 32)/0.001 = 16.002 \to 16$. \checkmark

\medskip
\textbf{Two paths.} Now let one quantity feed two others which recombine:
\[
  u = 2x, \qquad v = x^2, \qquad L = u + v .
\]
At $x = 3$: $u = 6$, $v = 9$, $L = 15$.

$\partial L/\partial u = 1$ and $\partial L/\partial v = 1$; $du/dx = 2$ and
$dv/dx = 2x = 6$. Influence arrives along both routes and the contributions
\emph{add}:
\[
  \frac{dL}{dx} = 1 \times 2 + 1 \times 6 = 8 .
\]
Check: $L = 2x + x^2$, so $dL/dx = 2 + 2x = 8$ at $x=3$. \checkmark

\medskip
\textbf{The two operations are the whole algorithm}: multiply along a path, add
across paths.
}

\formaldef{
If $L$ depends on $x$ through intermediate quantities $u_1, \ldots, u_m$, then
\[
  \frac{\partial L}{\partial x} = \sum_{j=1}^{m}
    \frac{\partial L}{\partial u_j}\,\frac{\partial u_j}{\partial x} .
\]
\textbf{Backpropagation} evaluates this from the output backwards. A
\emph{forward pass} computes and stores every intermediate value; a
\emph{backward pass} starts with $\partial L/\partial L = 1$ and walks toward the
inputs, multiplying by each local derivative and summing where paths meet.

Its efficiency is the point: one backward pass yields $\partial L/\partial w$ for
\emph{every} weight simultaneously, at roughly the cost of one forward pass ---
not one pass per weight.
}

\analogybreak{
The multiplication along a path is also the algorithm's weakness. Through $n$
layers the result is a product of $n$ local derivatives; if these are
systematically below 1 the product collapses toward zero, and if above 1 it
explodes. Neither is a flaw in the chain rule --- the algebra is exact --- but both
are real obstacles in floating-point arithmetic, and both are why activation
functions, initialisation schemes and gradient clipping receive so much
attention.
}

\whyhere{
This is how every model in this thesis is trained. It is also why the activation
choice in TimeGAN mattered so much: with sigmoid in the latent path, the local
derivative is at most $0.25$, so three layers multiply to at most $0.0156$ and the
learning signal reaching the early layers is negligible. Switching to $\tanh$,
whose derivative reaches $1$ at the origin, is what moved the autoencoder's $R^2$
from roughly zero to $+0.995$. The gradient-norm clip of $1.0$ used everywhere in
this project is a direct guard against the exploding case.
}

%% ─────────────────────────────────────────────────────────────────────────────
\section{The Jacobian}
\label{sec:jacobian}

\situation{
A machine takes two dials and drives three gauges. Turning either dial moves
potentially all three gauges. To describe the machine's responsiveness you need
one number for each dial--gauge pairing: two dials times three gauges is six
numbers, naturally laid out in a grid.
}

\problem{
The gradient of \S\ref{sec:gradient} was defined for functions returning a single
number. A neural network layer returns many numbers at once, so its
responsiveness is not one list but a whole table of them --- and the notation must
keep track of which entry means what.
}

\workedarith{
Take $F(x,y) = \bigl(x + y,\; xy,\; x^2\bigr)$: two inputs, three outputs.

Write the outputs as $F_1 = x+y$, $F_2 = xy$, $F_3 = x^2$ and differentiate each
with respect to each input:
\[
\begin{aligned}
  \partial F_1/\partial x &= 1, &\qquad \partial F_1/\partial y &= 1, \\
  \partial F_2/\partial x &= y, &\qquad \partial F_2/\partial y &= x, \\
  \partial F_3/\partial x &= 2x, &\qquad \partial F_3/\partial y &= 0 .
\end{aligned}
\]
Arrange as a grid, rows indexed by output and columns by input, evaluated at
$(x,y) = (2,3)$:
\[
  J = \begin{bmatrix} 1 & 1 \\ 3 & 2 \\ 4 & 0 \end{bmatrix} .
\]

\textbf{Reading it.} Row 2, column 1 is $3$: nudging $x$ moves $F_2$ at rate $3$.
Check: $F_2 = xy$, and at $(2.001, 3)$ it is $6.003$, up from $6$, so the rate is
$3$. \checkmark

Row 3, column 2 is $0$: $y$ has no effect on $F_3 = x^2$ at all. \checkmark
}

\formaldef{
For $F : \R^n \to \R^m$ with components $F_1, \ldots, F_m$, the
\textbf{Jacobian} is the $m \times n$ grid
\[
  J = \begin{bmatrix}
    \partial F_1/\partial x_1 & \cdots & \partial F_1/\partial x_n \\
    \vdots & \ddots & \vdots \\
    \partial F_m/\partial x_1 & \cdots & \partial F_m/\partial x_n
  \end{bmatrix} ,
  \qquad J_{ij} = \frac{\partial F_i}{\partial x_j} .
\]
When $m = 1$ the Jacobian has a single row and is the gradient. The chain rule for
composed multi-output functions is the product of their Jacobians, in the
grid-multiplication sense defined in \S\ref{sec:matmul}.
}

\analogybreak{
Jacobians are rarely built explicitly in practice. For a layer with $1{,}000$
inputs and $1{,}000$ outputs the grid holds a million entries, and a network has
many such layers. Automatic differentiation instead computes the \emph{product} of
a Jacobian with a given list of numbers, which needs far less memory and is all
backpropagation ever actually requires.
}

\whyhere{
Every layer of every network here has a Jacobian, and backpropagation is the
chained product of them. The critic's gradient in WGAN-GP is a one-row Jacobian ---
the critic outputs a single score --- which is why $\nabla D$ is a gradient and its
Euclidean length is the quantity the penalty constrains.
}

%% ─────────────────────────────────────────────────────────────────────────────
\section{The Hessian}
\label{sec:hessian}

\situation{
A valley may be a narrow gorge or a broad shallow bowl. Both have a flat floor,
so the gradient is zero in each. They demand completely different step sizes: the
gorge punishes a long stride by throwing you up the opposite wall, while the bowl
makes short strides a waste of time.
}

\problem{
The second-derivative test of \S\ref{sec:secondderiv} distinguished peaks from
valleys with one number. With many inputs there is curvature in every direction
and, worse, cross-terms describing how the slope in one direction changes as you
move in another. One number cannot carry that.
}

\workedarith{
Take $f(x,y) = x^2 + 3y^2$. First partials:
\[
  \frac{\partial f}{\partial x} = 2x , \qquad \frac{\partial f}{\partial y} = 6y .
\]
Both vanish only at $(0,0)$, so that is the sole candidate turning point.

Now differentiate each partial with respect to each variable:
\[
  \frac{\partial^2 f}{\partial x^2} = 2, \quad
  \frac{\partial^2 f}{\partial y^2} = 6, \quad
  \frac{\partial^2 f}{\partial x \partial y} = 0, \quad
  \frac{\partial^2 f}{\partial y \partial x} = 0 .
\]
As a grid:
\[
  H = \begin{bmatrix} 2 & 0 \\ 0 & 6 \end{bmatrix} .
\]

\textbf{Reading it.} Both diagonal entries are positive, so the surface curves
upward in both coordinate directions: $(0,0)$ is a minimum. The $y$ direction
curves three times as sharply as $x$, so the bowl is an elongated ellipse, steep
across and shallow along.

\textbf{Verify.} $f(1,0) = 1$ while $f(0,1) = 3$. The same unit step costs three
times as much in $y$. \checkmark

\medskip
\textbf{A saddle.} For $g(x,y) = x^2 - y^2$ the gradient $[2x, -2y]$ also vanishes
only at the origin, and
\[
  H = \begin{bmatrix} 2 & 0 \\ 0 & -2 \end{bmatrix} .
\]
One positive, one negative: a minimum along $x$, a maximum along $y$. Neither peak
nor valley but a saddle, and $g(1,0) = 1$ against $g(0,1) = -1$ confirms it. Such
points are everywhere in high dimensions --- with many directions available, it is
far more likely that some curve up and others down than that all agree.
}

\formaldef{
The \textbf{Hessian} of $f : \R^n \to \R$ is the $n \times n$ grid of second
partial derivatives,
\[
  H_{ij} = \frac{\partial^2 f}{\partial x_i \partial x_j} .
\]
For twice-continuously-differentiable $f$ the order of differentiation does not
matter, so $H_{ij} = H_{ji}$ and the grid is symmetric.

At a point where $\nabla f = \mathbf{0}$: if every direction curves upward the
point is a local minimum; if every direction curves downward, a local maximum; if
directions disagree, a saddle. Chapter~\ref{ch:linalg} makes ``every direction''
precise via eigenvalues (\S\ref{sec:eigen}).
}

\analogybreak{
The Hessian is the right diagnostic and the wrong computation. For a network with
$100{,}000$ weights it holds $10^{10}$ entries --- tens of gigabytes --- and must be
recomputed at every step. No training procedure in this thesis forms one.
Optimisers instead approximate curvature cheaply, which is exactly what Adam's
per-parameter scaling amounts to (Chapter~10 (Optimisation)).
}

\whyhere{
Curvature explains why a single global learning rate struggles: in a valley shaped
like the $[2,0;0,6]$ example, the step that suits one direction is wrong for the
other. Adam --- the optimiser used for every model in this project --- addresses
precisely this by maintaining a separate effective step size per parameter.
}
